\section{Introduction}\label{sec:introduction}

Let $X$ be a smooth projective connected curve over $\mathbb C$, and let $G$
be a connected simple complex algebraic group. Its Langlands dual
$G^\vee$ is defined by the dual of the full root datum, so the global form
of $G$ is part of the data. Write $\operatorname{Bun}_G(X)$ for the stack
of principal $G$-bundles. The adjoint determinant line on this stack is
\[
  L_G|_E=\det R\Gamma(X,\mathfrak g_E).
\]
Thus the determinant of $H^1$ occurs with exponent $-1$. If $h^\vee$ is
the dual Coxeter number, set
\begin{equation}\label{intro:level}
  D_c(\operatorname{Bun}_G(X))
  :=D\!\left(L_G^{(c-h^\vee)/(2h^\vee)}\right).
\end{equation}
Here $D(L^s)$ is the cocomplete derived DG category of all twisted
$D$-modules for the indicated complex power of the line. The power specifies
a sheaf of twisted differential operators and does not require a root of
$L$. Precise descent and sign conventions are given in
Section~\ref{sec:foundations}.

Let $r$ be the lacing number of $G$: it is $1$, $2$, or $3$ according to
the ratio of the squared root lengths. Fix dual pinnings and a theta
characteristic on $X$. These choices fix the residue characters and the
normalizations of localization used below.

\begin{theorem}\label{thm:main}
For every $c\in\mathbb C\setminus\mathbb Q$ there is an equivalence of
presentable $\mathbb C$-linear DG categories
\begin{equation}\label{intro:equivalence}
  D_c(\operatorname{Bun}_G(X))
    \simeq
  D_{-1/(rc)}(\operatorname{Bun}_{G^\vee}(X)).
\end{equation}
The construction uses the given global forms and all connected components.
It is characterized by the localization--Whittaker comparison of
Section~\ref{sec:equivalence}, including its compatibility with families of
marked points and their collisions. It also intertwines central torsor actions
with the corresponding transgression local systems.
\end{theorem}

The normalization \eqref{intro:level} is the shifted-level convention of
\cite[Section~1.1.5, Conjecture~1]{BogdanovaBetti}. In particular,
Theorem~\ref{thm:main} gives its unramified de Rham equivalence in the
irrational range. The localization comparison and the central compatibility
are part of the construction: they identify the functor as the quantum
Langlands correspondence. Their definitions and formulas appear with the
functors they concern in Section~\ref{sec:equivalence}.

\subsection{Context and prior work}

Geometric Langlands relates sheaf-theoretic categories on a bundle stack to
categories built from local systems for the dual group. Its quantum form
compares twisted $D$-module categories at inverse shifted levels. The
gauge-theoretic interpretation through electric--magnetic duality, developed
by Kapustin and Witten~\cite{KapustinWitten}, explains the appearance of the
dual root datum and reciprocal parameter. An early formulation of the
quantum deformation, attributed there to Drinfeld, appears in Stoyanovsky's
inverse-parameter conjecture for twisted $D$-module
categories~\cite[Section~1.3, Conjecture~2]{Stoyanovsky}. The formulation in
terms of local categories, localization, and Whittaker coefficients is
developed in Gaitsgory's quantum Langlands program~\cite{GaitsgoryQGL}.

Several established results supply the local and classical parts of the
argument. Feigin--Frenkel duality identifies principal $W$-algebras at inverse
shifted levels~\cite{FeiginFrenkel}. Arakawa--Frenkel identify the reductions of Weyl modules and
their coweight twists, including concentration and simplicity at irrational
level~\cite{ArakawaFrenkel}. Raskin's categorical form of
Drinfeld--Sokolov reduction gives a precise renormalized $W$-module category
and generalized vacuum objects~\cite{RaskinW}. The fundamental local
equivalences of Campbell--Dhillon--Raskin give the required irrational-level
point categories for the actual global forms~\cite{CDR}. Gaitsgory's
local/global Whittaker comparison supplies the bridge between the affine
Grassmannian and global pole models, also over multipoint parameter
schemes~\cite{GaitsgoryWhit}.

On the global side, the proof of ordinary geometric Langlands provides the
spectral description of the tempered category and its enhanced
Whittaker--Poincar\'e comparison~\cite{GLCI,GLCII,GLCIII,GLCIV,GLCV}.
The arbitrary-level localization and unipotent-period constructions of
\cite{GLCII} supply further inputs. In
Sections~\ref{sec:localization}--\ref{sec:ran} we construct the period
comparison with the finite levels, arbitrary Grassmannian tests, and
moving-point operations needed below.
Drinfeld--Gaitsgory's finiteness and
constant-term results and Gaitsgory's miraculous duality furnish the
categorical and geometric framework for our twisted
duality~\cite{DrinfeldGaitsgoryFiniteness,DrinfeldGaitsgoryCT,GaitsgoryStrange}.
Lin's comparison of Poincar\'e series with miraculous duality supplies the
ordinary-level geometric maps that we deform~\cite{Lin}.

Recent work of Bogdanova constructs a quantum Langlands functor in the
Betti setting and specifies its central compatibility through
2-Fourier--Mukai duality~\cite{BogdanovaBetti}. Her work on quantum
Whittaker nonvanishing also makes explicit the finite local averaging
diagram used in our argument~\cite{BogdanovaWhit}. The two further
mechanisms developed here are a
weight argument controlling filtered deformation of a global kernel,
and an extension argument turning local pairings into coherent
comparisons across collisions.

\subsection{The two main comparisons}

Localization takes representations of loop Lie algebras to twisted
$D$-modules on the bundle stack and is a quotient functor. We first
compare the two sides on localized objects, using local pairings and a
global period identity. The quotient property then produces a functor
between the bundle categories; two global Whittaker identities provide
its left and right inverses. We explain the two main inputs to this
construction.

The first concerns global Whittaker coefficients. For a group $H$ and a
determinant exponent $s$, let $R_s$ denote the coefficient functor with
arbitrary moving sets of allowed poles. Its target remembers the
nondegenerate unipotent character, as well as the pole positions. Pairing
with the opposite character and applying the twisted pseudo-identity
produces a Poincar\'e functor $J_s$. We construct a comparison and prove
that it is an isomorphism:
\begin{equation}\label{intro:plancherel}
                 J_sR_s\longrightarrow\operatorname{Id}
                 \qquad(s\notin\mathbb Q).
\end{equation}
This comparison is defined on the entire twisted $D$-module category.
Its ordinary-level counterpart has an anti-tempered remainder. A fixed
finite Whittaker test annihilates that remainder and is conservative at
irrational exponent. It therefore suffices to prove that the tested
kernel cone vanishes after deformation.

Finite pole bounds and finite diagrams of marked-point maps give finite
approximations to this cone. We express them by ordinary geometric
complexes with a universal Kummer local system of monodromy $q$;
Fourier transform in auxiliary affine lines then introduces the
exponential characters. Quasi-unipotent monodromy makes cohomology ranks
and transition-map ranks constant away from roots of unity. The passage
to the colimit requires a further argument. Compatible logarithm
variations and a uniform upper weight bound imply that each cohomology
generator in the completion at $q=1$ is killed by a finite transition.
Constancy of transition ranks then gives the same vanishing at every
non-torsion $q$, hence at every irrational exponent. Ordinary geometric
Langlands supplies the cohomological bound used to make the weights
uniform. Section~\ref{sec:deformation} proves the deformation principles;
Section~\ref{sec:plancherel} applies them to the tested kernel cone.

The second comparison concerns local categories over families of points.
For $G$ and $G^\vee$, convolution followed by Drinfeld--Sokolov
reduction pairs a positive-loop-equivariant Kac--Moody module with a
Whittaker object on its affine Grassmannian. At a single point,
inverse-level $W$-algebra duality identifies reductions with the
representation and modification labels exchanged. To use these pairings
globally, one needs maps over the
whole parameter scheme, including collision diagonals, together with
their insertion and factorization compatibilities. We construct the two
pairings with a common $W$-module target and continuous relative right
adjoints. On each configuration stratum, the cohomology of the
semisimple group makes the adjunction units isomorphisms through degree
two on heart objects. Maps between heart objects can therefore be
recovered from their images in the common target. Restriction across a
collision divisor has amplitude $[0,1]$; the recovered maps identify its
boundary epimorphisms, extending the equivalence of hearts across the
divisor. Realization and the adjunctions give the equivalence of DG
categories. A separate extension argument, using translation of the
marked points and the horizontal vacuum, identifies the pairings with
both slots inserted. Sections~\ref{sec:loop}--\ref{sec:local} carry out
these steps.

\subsection{Organization and the global period comparison}

The category-theoretic setup begins in Section~\ref{sec:foundations}.
Irrational inertia weights eliminate sufficiently unstable
Harder--Narasimhan strata. This reduces the relevant bundle categories
to quasi-compact opens and allows ordinary categorical duals to be used.
Section~\ref{sec:pseudoidentity} proves twisted pseudo-identity using
compactifications, adjunctions, and the Levi boundary of the diagonal.
Section~\ref{sec:whittaker} fixes the Whittaker models and their dual
pairings.

After the two main comparisons, Sections~\ref{sec:localization}--\ref{sec:ran}
identify unipotent periods with integrated reduction complexes
(Theorem~\ref{loc:period-comparison}). The
proof keeps the determinant lines, finite-level shifts, and moving-point
operators throughout. Ran integration contracts the added point
variables and turns the local pairing comparison into a global period
identity.

Section~\ref{sec:equivalence} assembles the equivalence through the
localization quotient, using the period identity and the two instances
of \eqref{intro:plancherel}. It also derives the exchange of components
and central characters and the coherent action of central torsors.
